\section{The original strip has linear elimination depth; its leading native covariance loss is an original projection angle}\label{the-original-strip-has-linear-elimination-depth-its-leading-native-covariance-loss-is-an-original-projection-angle}

21 September 2026. Independent continuation of CGR1--29. The full original conductor, invariant image, Gamma source and quotient minimum are retained. The results below concern the actual row space of CG29--30 at all four original cutoffs. They quantify the source contribution and locate the remaining leading covariance in an exact projection. They do not assign the remaining projection determinant a value.

\subsection{1. The exact conductor elimination has depth at most 10k}\label{the-exact-conductor-elimination-has-depth-at-most-10k}

Use the entire OCF4 conductor \[
A=\sum_{r,s=0}^{8}a_{rs}e_{rs}^{(8)},\qquad
(r_0,s_0)=\max_{\mathrm{lex}}\{(r,s):a_{rs}\ne0\},\qquad
a_*=a_{r_0s_0},\qquad A_1=\sum_{r,s}|a_{rs}|.
\tag{CGS1}
\] The subscripts \(r_0,s_0\) here denote the actual leading monomial, not the earlier residual-rank symbol. Set \[
B_A^{\mathrm{step}}=\frac{A_1-|a_*|}{|a_*|},\qquad
C_A^{\mathrm{path}}=1+B_A^{\mathrm{step}}=\frac{A_1}{|a_*|}\ge1.
\tag{CGS2}
\] Give the original coefficient position \((a,b)\in\{0,\ldots,k\}^2\) the integer height \(h(a,b)=9a+b\). OCF6 replaces a pivot coefficient at \((r_0+i,s_0+j)\) by the entire sum \[
-\sum_{(r,s)\ne(r_0,s_0)}\frac{a_{rs}}{a_*}
   e_{r+i,s+j}^{(k)}.
\tag{CGS3}
\] Every nonzero term has strictly smaller height. If \(r=r_0\), its decrease is \(s_0-s\ge1\). If \(r<r_0\), the decrease is \(9(r_0-r)+(s_0-s)\ge9-8=1\). Every destination remains in the original coefficient box, since \(0\le i,j\le k-8\) and \(0\le r,s\le8\). Thus a path of successively replaced coefficients has length at most \(10k\); coefficients are never removed from the original box to obtain this bound.

For completeness, expand the reduction of one basis column as the sum over all its substitution paths, stopping each path when it reaches the actual strip. At a substituted node the sum of absolute edge weights is \(B_A^{\mathrm{step}}\). The sum of absolute weights of all terminal paths of a specified length \(j\) is consequently at most \((B_A^{\mathrm{step}})^j\), by induction on the successive levels. The triangle inequality therefore bounds the total output column norm by \(\sum_{j=0}^{10k}(B_A^{\mathrm{step}})^j\), which is at most \((1+B_A^{\mathrm{step}})^{10k}\). The latter assertion follows by expanding the binomial and using \(\binom{10k}{j}\ge1\). When \(B_A^{\mathrm{step}}=0\), the same expression is one. The exact strip map therefore obeys \[
\boxed{\|N_k\|_{1\to1}\le(C_A^{\mathrm{path}})^{10k},\qquad
\|N_k\|_{2\to2}\le\sqrt q(C_A^{\mathrm{path}})^{10k}.}
\tag{CGS4}
\] This computes a quantitative bound for the existing OCF elimination; the order and coefficients of that elimination have not been changed. The second inequality uses \(\|u\|_2\le\|u\|_1\) on the output and \(\|f\|_1\le\sqrt q\|f\|_2\) on its original input.

In the identical strip positions, \(N_kz=z\) and \(N_kM_A=0\). Thus the complete conductor image has the following proved angle from every original strip coefficient vector: \[
\boxed{\operatorname{dist}_{\ell^2}(z,\operatorname{im}M_A)
\ge\frac{\|z\|_2}{\sqrt q(C_A^{\mathrm{path}})^{10k}}.}
\tag{CGS5}
\] Apply \(N_k\) to \(z-M_Au\), use CGS4, and take the infimum over all original conductor coefficients \(u\) to prove the formula. Consequently CGR15 remains valid with the path-based lower constant \[
l_{U,N}^{\mathrm{path}}=
\left[M_D^2B_{\mathrm{val}}q
          (C_A^{\mathrm{path}})^{20k}\right]^{-1}.
\tag{CGS6}
\] The \(B_{\mathrm{val}}\) in this statement is still the complete \(q\)-root interpolation constant of CGR13. The original source projection is still present; CGS5 does not erase that factor. One may take the larger of CGS6 and the prior lower constant at each finite cutoff. The new path bound improves the growth exponent as \(k\) increases; it need not be the smaller numerical bound for the first few allowed values of \(k\).

\subsection{2. The sparse strip has a smaller interpolation constant before projection}\label{the-sparse-strip-has-a-smaller-interpolation-constant-before-projection}

There are exactly \[
n_E=|E_k^{\mathrm{strip}}|=q-q'=16k-48
\tag{CGS7}
\] original roots in the strip. Their physical coordinates are \[
y_{ab}=(\beta_{ab}-c)/i=(2b-k)\gamma-i(2a-k)\delta,
\quad (a,b)\in E_k^{\mathrm{strip}},\quad c=k/2.
\tag{CGS8}
\] Every one has modulus at most \(R_y=k\sqrt{\delta^2+\gamma^2}\). Every pair has separation at least \(2\delta\), because \(0<\delta<1/2\), \(\gamma>2\), and distinct pairs differ in at least one coordinate. In the original Gamma norm define its Lagrange polynomials by \[
L_\alpha(y)=\prod_{\beta\in E_k^{\mathrm{strip}},\ \beta\ne\alpha}
\frac{y-y_\beta}{y_\alpha-y_\beta}.
\tag{CGS9}
\] These interpolate precisely the strip roots. No lower-lattice root is inserted in this interpolation and none is dropped from the later source projection.

Successive application of CG11 to every displayed factor gives \[
\|L_\alpha\|_\sigma^2
\le M_\sigma\left(\frac{2n_E+R_y}{2\delta}\right)^{2(n_E-1)},
\qquad
B_E=n_EM_\sigma
\left(\frac{2n_E+R_y}{2\delta}\right)^{2(n_E-1)}.
\tag{CGS10}
\] Indeed the degree before any multiplication is at most \(n_E-2\), so CG11 bounds that multiplication norm by \(2n_E+R_y\); each denominator has modulus at least \(2\delta\), and the norm of the constant polynomial is exactly \(\sqrt{M_\sigma}\).

Let \(E_{E,N}\) evaluate the complete upper Gamma frame of degrees \(0,\ldots,N\) at these actual roots. Since \(n_E-1\le q-1\le N\), the polynomials CGS9 form a right inverse of this evaluation map at every original cutoff. Its squared operator norm is at most the sum of the squared column norms, namely \(B_E\). For a row \(z\), applying that right inverse to \(zE_{E,N}\) gives \[
\|zE_{E,N}\|_2^2\ge B_E^{-1}\|z\|_2^2.
\tag{CGS11}
\] This proves the lower bound without a Vandermonde determinant estimate.

The actual rows in \(Z\) are supported on this strip and annihilate all polynomials of degree below \(v\). Their first \(v\) Gamma functional coefficients consequently vanish exactly. Thus removing those zero columns leaves the identical row norm, and CGR7 satisfies \[
\boxed{B_E^{-1}R_Z\preceq B_FB_F^*\preceq V_NR_Z.}
\tag{CGS12}
\] The right inequality is CGR10. The left inequality uses exactly the unchanged sparse rows of the invariant image. It is not a statement about the covariance after projecting away the lower root evaluations.

\subsection{3. The complete unprojected native covariance has only o(kq) determinant cost}\label{the-complete-unprojected-native-covariance-has-only-okq-determinant-cost}

Retain \(A=A_{D,1}\), \(D=N-v\), and define \[
W_N=B_FA^{-1},\qquad
\widetilde C_{R,N}=W_NW_N^*,\qquad
C_{R,N}^U=W_NP_IW_N^*.
\tag{CGS13}
\] Every matrix is in the original lower-center Gamma frame; CGR9 proves the last equality with the original orthogonal projection \(P_I\) onto \(\chi'\mathcal P'_{L+g}\). The forward WCF bound \(\|A\|\le M_D\) and CGS12 prove \[
\widetilde C_{R,N}\succeq
           (M_D^2B_E)^{-1}R_Z.
\tag{CGS14}
\] For each \(j\le r\), the same exterior argument as CGR11, now without a projection, gives \[
\prod_{i=1}^j\widetilde\lambda_i
\le V_N^jM_D^{2(D+1-j)}e^{2J_D},
\tag{CGS15}
\] where the ordered eigenvalues are those relative to \(R_Z\). All factors come from the original operator and the exact strip rows.

Set \[
\begin{split}
K_{-,N}&=r\log^+(M_D^2B_E),\\
K_{+,N}&=r\log^+V_N+2(D+1)\log M_D+2J_D,\\
K_N&=\max\{K_{-,N},K_{+,N}\}.
\end{split}\tag{CGS16}
\] Taking the determinant in CGS14 bounds its logarithm below by \(-K_{-,N}\). Taking the largest partial sum of the ordered logarithms in CGS15 bounds its positive logarithmic part by \(K_{+,N}\), hence its whole logarithm above by that number. We have therefore evaluated the size of the complete source contribution: \[
\boxed{\left|\log\frac{\det\widetilde C_{R,N}}{\det R_Z}\right|
\le K_N=O_A(q\log(q+2))=o(kq).}
\tag{CGS17}
\] To check the stated order, \(r=O(k)\), \(n_E=O(k)\), \(R_y=O_A(k)\), \(\log B_E=O_A(k\log(k+2))\), \(\log V_N=O_A(k\log(q+2))\), and \(\log M_D=O_A(\log(q+2))\), \(J_D=O_A(q)\) are the proved bounds of the cited formulas. The endpoints satisfy \(D\le2q\). Thus both finite quantities in CGS16 are \(O_A(q\log(q+2))\); since \(q=(k+1)^2\), division by \(kq\) tends to zero.

\subsection{4. An exact principal-angle formula for the uncomputed native loss}\label{an-exact-principal-angle-formula-for-the-uncomputed-native-loss}

The following maps give the geometry of the remaining term, with their domains and codomains retained: \[
\mathbb C^r\xrightarrow{\ T_N=W_N^*\widetilde C_{R,N}^{-1/2}\ }
\mathcal P'_D\xrightarrow{\ P_I\ }
I_N=\chi'\mathcal P'_{L+g}\subset\mathcal P'_D.
\tag{CGS18}
\] The middle space has the original center-\(c'\) Gamma inner product and its orthonormal coefficient frame. The first map is an isometry: \(T_N^*T_N=I_r\) follows directly from CGS13. The second map is the original orthogonal projection and uses every lower-root condition. Consequently \[
S_N=T_N^*P_IT_N
=\widetilde C_{R,N}^{-1/2}C_{R,N}^U
                   \widetilde C_{R,N}^{-1/2}
\quad\hbox{satisfies}\quad 0<S_N\preceq I_r.
\tag{CGS19}
\] Strict positivity follows because the original quotient observation has full row rank: if a combination of its rows has zero norm on \(I_N\), it vanishes on the low target polynomials too, contrary to the full row rank of \(A_R\). The eigenvalues of \(S_N\) are the squares of the singular values of \(P_IT_N\), or equivalently the squared cosines of the principal angles between \(\operatorname{im}T_N\) and \(I_N\). This equivalence follows from the defining Gram \((P_IT_N)^*(P_IT_N)=S_N\); it introduces no assumed angle bound.

Taking determinants in CGS19 and applying CGS17 proves the finite quantitative receiver \[
\boxed{\left|\log\frac{\det C_{R,N}^U}{\det R_Z}
                 -\log\det S_N\right|\le K_N=o_A(kq).}
\tag{CGS20}
\] Thus the entire leading native covariance loss is the sum of logarithms of these original squared projection angles. Every contribution from the unprojected source, including its actual invariant coefficients and the full conductor inverse, has been bounded at the required scale.

For the two high endpoints, put \(D_{d,N}=\log(\det C_{R,N}^d/\det R_Z)\), with the unchanged outer-ideal covariance of CG31. The earlier complete graph calculation then gives \[
\left|\mathcal G_K-
\sum_{N\in\{2q-1,2q\}}(D_{d,N}-\log\det S_N)\right|
\le B_{\mathrm{CG32}}+K_{2q-1}+K_{2q}=o_A(kq),
\tag{CGS21}
\] where \(B_{\mathrm{CG32}}\) is exactly the finite right side of CG32. To verify the sign, CG31 defines \(\Delta_{R,N}=\log\det C_{R,N}^U-\log\det C_{R,N}^d\), and CG32 gives \(\mathcal G_K=-\sum_N\Delta_{R,N}\) up to its displayed error. Insert CGS20 into this identity.

Every statement CGS12--20 also applies to the specified pole-canceling row subspace \(\mathcal W_k^0\) of CGR23, using the same restricted coefficient form and its actual row dimension. On that subspace CGR26 also bounds the unprojected covariance, because its proof bounds the full Gamma functional before using the projection. Hence \[
(M_D^2B_E)^{-1}R_Z|_{\mathcal W_k^0}
\preceq\widetilde C_{R,N}|_{\mathcal W_k^0}
\preceq e^{\varepsilon_{k,N}}R_Z|_{\mathcal W_k^0}.
\tag{CGS22}
\] Both logarithmic quadratic-form allowances are \(o(q)\), with all constants given in CGS10 and CGR26. Combining the subspace version of CGS21 with CGR28 retains its additional, already proved \(O_A(q\sqrt{k}\log^2(q+2))\) error. The unresolved leading quantity is now precisely the relative outer covariance versus these original projection angles. No value or sign for that quantity is asserted by the present bound.

\subsection{Sources and coverage}\label{sources-and-coverage}

OCF1--26, \emph{An exact conductor frame and a fixed invariant-generator recursion}, was read completely. CGS1--5 use OCF4--8 with the identical coefficients, pivot and original strip. The \href{https://github.com/KokunoYumeto/zeta-function-research-reader/blob/5992ad50c0f2a06921f1fcaded7b4e38097c0739/workbenches/splitzero-tandem/continuations/20260920-original-kernel-mixed-determinants/providers/ORIGINAL_CONDUCTOR_STRIP_FRAME.tex\#L27}{complete public proof file} has a byte-verification receipt in the parent source ledger.

KF1--55, CG1--32 and CGR1--29 were read completely. The source files are \nolinkurl{ACTUAL_KERNEL_FILTRATION_AND_GRAPH.md}, \nolinkurl{GRAPH_METRIC_DERIVATION.md} and \nolinkurl{INVARIANT_COVARIANCE_POLE_FILTRATION.md}, respectively. Their complete sources must accompany public use when a pinned proof-file link has not yet been verified. Repeated equation prefixes in other documents do not identify these sources.

The WCF forward and exterior bounds are WCF4--13 and WCF19a--f in the \href{https://github.com/KokunoYumeto/zeta-function-research-reader/blob/52274804051e325c8a6ef69c17b3958cc2caba0a/workbenches/splitzero-tandem/continuations/20260920-gamma-and-finite-derivative/providers/WEIGHTED_CONDUCTOR_FORWARD.tex}{complete pinned proof file}. The cited portions were read directly, and the public source was checked byte for byte against the local source used here. CG11 derives its Gamma factor estimate from the Meixner--Pollaczek recurrence and generating function in NIST DLMF \href{https://dlmf.nist.gov/18.22.E8}{18.22.8} and \href{https://dlmf.nist.gov/18.23.E7}{18.23.7}, by T. H. Koornwinder, R. Wong, R. Koekoek, R. F. Swarttouw and W. P. Reinhardt. The original NIST TeX encodings are retained and read in \nolinkurl{graph_primary_sources/}. The path-depth proof, all interpolation inequalities and the exact projection-angle map used here are proved in this file.

The independent \nolinkurl{check_invariant_strip.py} passed 282 exact checks. They include all 81 possible leading-monomial height cases; the complete original decreasing-lex elimination at eight auxiliary conductor fixtures; the sparse Gamma interpolation estimate; and the exact retained projection-angle determinant. A negative control verifies that the projection determinant cannot be set equal to one. The finite fixtures support the algebraic identities and do not evaluate the remaining original growing-degree determinant.

The exact earlier proof files are \href{https://github.com/KokunoYumeto/zeta-function-research-reader/blob/0fd693c844aad9117e30ec447c59864a28af10f3/workbenches/splitzero-tandem/continuations/20260921-native-spectral-real-pair/003/ACTUAL_KERNEL_FILTRATION_AND_GRAPH.tex\#L12}{KF1--55}, \href{https://github.com/KokunoYumeto/zeta-function-research-reader/blob/0fd693c844aad9117e30ec447c59864a28af10f3/workbenches/splitzero-tandem/continuations/20260921-native-spectral-real-pair/003/COMPLETE_CONDUCTOR_GRAPH_METRIC.tex}{CG1--32}, and \href{https://github.com/KokunoYumeto/zeta-function-research-reader/blob/0fd693c844aad9117e30ec447c59864a28af10f3/workbenches/splitzero-tandem/continuations/20260921-native-spectral-real-pair/003/COMPLETE_CONDUCTOR_ROOT_GEOMETRY.tex\#L38}{OR1--25}. CGR and CGS are reproduced completely in this edition.
